% Hough 变换教学例：三点 (1,3),(2,2),(3,1) 位于直线 x+y=4 上；右图为
% rho=x cos(theta)+y sin(theta) 的三条正弦曲线，交于 (45 度, 2.83)。宽约 104mm。
\begin{tikzpicture}[x=1mm,y=1mm,font=\figurefont\small,
  ax/.style={-{Stealth[length=1.8mm]},FigureStroke,line width=.6pt},
  lab/.style={font=\figurefont\footnotesize}]
  % image plane, 1 unit = 8mm
  \begin{scope}[x=8mm,y=8mm]
    \draw[ax] (0,0)--(4.4,0) node[right] {$x$};
    \draw[ax] (0,0)--(0,4.4) node[above] {$y$};
    \foreach \t in {1,2,3,4}{\draw[FigureStroke] (\t,0)--(\t,-.08) node[below,lab] {\t};
      \draw[FigureStroke] (0,\t)--(-.08,\t) node[left,lab] {\t};}
    \draw[FigureBlue,line width=.9pt] (0.3,3.7)--(3.7,0.3);
    \draw[-{Stealth[length=1.8mm]},FigureYellow,line width=.9pt] (0,0)--(2,2);
    \node[lab,anchor=north west] at (1.05,1.0) {$\rho$};
    \draw[FigureYellow,line width=.6pt] (0.7,0) arc[start angle=0,end angle=45,radius=0.7];
    \node[lab] at (0.95,0.32) {$\theta$};
    \foreach \x/\y in {1/3,2/2,3/1}{\fill[FigureRed] (\x,\y) circle (0.9mm);}
    \node at (2.2,4.9) {图像平面};
  \end{scope}
  % parameter plane: theta 0..180 deg -> 0..60mm, rho -3.5..4.5 -> scale 6mm
  \begin{scope}[xshift=44mm,yshift=22mm]
    % theta axis at rho=-3.5 (bottom), rho=0 as grid line
    \draw[FigureGrid,line width=.5pt] (0,0)--(50,0);
    \draw[ax] (0,-21)--(54,-21) node[right] {$\theta$};
    \draw[ax] (0,-21)--(0,24) node[above] {$\rho$};
    \foreach \d in {0,45,90,135,180}{\draw[FigureStroke] ({\d/3.6},-21)--({\d/3.6},-22) node[below,lab] {$\d^\circ$};}
    \foreach \r in {-3,0,3}{\draw[FigureStroke] (0,{6*\r})--(-1,{6*\r}) node[left,lab] {$\r$};}
    \tikzset{houghcurve1/.style={FigureBlue},houghcurve2/.style={FigureStroke,dashed},houghcurve3/.style={FigureRed}}
    \foreach \px/\py/\c in {1/3/houghcurve1,2/2/houghcurve2,3/1/houghcurve3}{
      \draw[\c,line width=.9pt,domain=0:180,samples=60,smooth]
        plot ({\x/3.6},{6*(\px*cos(\x)+\py*sin(\x))});
    }
    \node[lab,anchor=west,text=FigureBlue] at (50.5,-6) {$(1,3)$};
    \node[lab,anchor=west,text=FigureStroke] at (50.5,-12) {$(2,2)$};
    \node[lab,anchor=west,text=FigureRed] at (50.5,-18) {$(3,1)$};
    \fill[FigureStroke] (12.5,{6*2.8284}) circle (1mm);
    \draw[FigureMuted,dashed,line width=.5pt] (12.5,-21)--(12.5,{6*2.8284});
    \node[lab,anchor=south] at (12.5,{6*2.8284+1.5}) {$(45^\circ,\,2.83)$};
    \node at (25,30) {参数平面};
  \end{scope}
\end{tikzpicture}
